\bookchapter{Κεφάλαιο 54: Έργο 1 — Το βιβλίο εξόδων $\mu$ου (τοπική έκδοση)}

\chapterintro{Καλώς ήρθες στο πρώτο σου «κανονικό» έργο. Αυτό το βιβλίο εξόδων, παρότι δεν συνδέεται στο διαδίκτυο, είναι πλήρες σε όλα: καταγράφει συναλλαγή, δείχνει την κίνηση, κάνει $\mu$ηνιαία στατιστικά, επιτρέπει διαγραφή και αλλαγή. Όταν το τελειώσεις, θα έχεις περπατήσει ολόκληρη τη διαδρο$\mu$ή «δεδο$\mu$ένα λογική διεπαφή». Μη βιάζεσαι να αντιγράψεις κώδικα — πρώτα σκέψου καθαρά, $\mu$ετά δράσε· αυτός είναι ο $\mu$όνι$\mu$ος κανόνας $\mu$ας.}

\booksection{Σκέψη απαιτήσεων: ρώτα πρώτα πέντε ερωτήσεις}

Πριν πας στη δράση, ξεκαθάρισε την ιδέα $\mu$ε τα «πέντε στοιχεία απαιτήσεων»:

\begin{center}
\begin{tabularx}{\linewidth}{|X|X|}
\hline
\textbf{Στοιχείο} & \textbf{Απάντηση} \tabularnewline
\hline
Who (για ποιον) & Για $\mu$ένα — και για οικογένεια και φίλους \tabularnewline
What (τι κάνει) & Καταγραφή εσόδων-εξόδων, προβολή κίνησης, $\mu$ηνιαία σύνοψη, στατιστικά ανά κατηγορία \tabularnewline
Why (γιατί) & Εξάσκηση + να καταλάβω πού πάνε τα λεφτά $\mu$ου \tabularnewline
How (πώς) & Kotlin + Compose + Room + MVVM, α$\mu$ιγώς τοπικό χωρίς backend \tabularnewline
Constraint (περιορισ$\mu$οί) & Ολοκλήρωση $\mu$έσα σε $\mu$ία εβδο$\mu$άδα· ελάχιστη υποστήριξη Android 8· έστω χωρίς δη$\mu$οσίευση, να γράφεται $\mu$ε πρότυπα δη$\mu$οσίευσης \tabularnewline
\hline
\end{tabularx}
\end{center}

\begin{center}\small\color{gray}Τα πέντε στοιχεία απαιτήσεων του βιβλίου εξόδων\end{center}

Μετά σπά’ το σε $\mu$ικρά βή$\mu$ατα «κάθε ένα χρήσι$\mu$ο από $\mu$όνο του»: (1) φτιάξε το στρώ$\mu$α δεδο$\mu$ένων Room (να αποθηκεύει $\mu$ια εγγραφή)· (2) σελίδα λίστας (να βλέπεις την κίνηση)· (3) σελίδα καταχώρησης (να προσθέτεις εγγραφές)· (4) διαγραφή και επεξεργασία· (5) κάρτα $\mu$ηνιαίας σύνοψης. Σε κάθε βή$\mu$α που ολοκληρώνεις, η εφαρ$\mu$ογή είναι ήδη κάτι «χρήσι$\mu$ο».

\booksection{Στρώ$\mu$α δεδο$\mu$ένων: να πάρουν οι «εγγραφές» σπίτι}

Ορίσε πρώτα πώς $\mu$οιάζει $\mu$ια εγγραφή — τα πεδία ας φτάνουν για τις ανάγκες, χωρίς απληστία:

\begin{lstlisting}[language=Kotlin]
@Entity(tableName = "records")
data class Record(
    @PrimaryKey(autoGenerate = true) val id: Long = 0,
    val type: String,        // "income" έσοδο | "expense" έξοδο
    val amount: Double,      // Ποσό (γιουάν)
    val category: String,    // Κατηγορία: φαγητό, $\mu$ετακινήσεις, $\mu$ισθός……
    val note: String = "",   // Ση$\mu$είωση (προαιρετικό)
    val date: String         // "2026-08-15"
)
\end{lstlisting}

Στο Dao πρόσθεσε δύο ερωτή$\mu$ατα «για στατιστικά» — αυτή είναι η ψυχή του βιβλίου εξόδων:

\begin{lstlisting}[language=Kotlin]
@Dao
interface RecordDao {
    @Query("SELECT * FROM records ORDER BY date DESC, id DESC")
    fun getAll(): Flow<List<Record>>

    @Insert
    suspend fun insert(record: Record)

    @Delete
    suspend fun delete(record: Record)

    @Query("SELECT COALESCE(SUM(amount),0) FROM records WHERE type='income' AND date LIKE :month || '%'")
    suspend fun monthIncome(month: String): Double   // month σε $\mu$ορφή "2026-08"

    @Query("SELECT COALESCE(SUM(amount),0) FROM records WHERE type='expense' AND date LIKE :month || '%'")
    suspend fun monthExpense(month: String): Double
}
\end{lstlisting}

\begin{tcbinfo}{Συ$\mu$βουλή: γιατί το ποσό είναι Double;}
Τα αυστηρά χρη$\mu$ατοοικονο$\mu$ικά προϊόντα χρησι$\mu$οποιούν BigDecimal ή «Int σε $\mu$ονάδες λεπτών (cents)» για να αποφύγουν σφάλ$\mu$ατα κινητής υποδιαστολής (το παλιό πρόβλη$\mu$α 0.1+0.2$\neq$0.3). Για έργο εξάσκησης το Double αρκεί· αν κάποτε φτιάξεις πραγ$\mu$ατικό χρη$\mu$ατοοικονο$\mu$ικό προϊόν, θυ$\mu$ήσου να αλλάξεις σε $\mu$ονάδα «λεπτά» — είναι παγίδα στην οποία έχουν πέσει α$\mu$έτρητοι προηγού$\mu$ενοι.

\end{tcbinfo}

\booksection{MVVM: διεπαφή και λογική «χωρίζουν σπίτια»}

Να αφήσεις τη διεπαφή να $\mu$ιλάει απευθείας $\mu$ε τη βάση; Θα τρέχει, αλλά ο κώδικας της διεπαφής γίνεται ατελείωτος και ακο$\mu$ψός. Η αρχιτεκτονική MVVM $\mu$οιράζει καθήκοντα: η View (διεπαφή) $\mu$όνο ζωγραφίζει, το ViewModel (κέντρο λογικής) διαχειρίζεται τα δεδο$\mu$ένα, το Model (Room) την αποθήκευση.

View (οθόνη Compose)

ViewModel (κατάσταση και λογική)

Model (στρώ$\mu$α δεδο$\mu$ένων Room)

\begin{center}\small\color{gray}Εικόνα 5: Η τριών επιπέδων δο$\mu$ή του MVVM — τα δεδο$\mu$ένα ρέουν προς $\mu$ία κατεύθυνση και κάθε επίπεδο $\mu$ένει στα σύνορά του\end{center}

\begin{lstlisting}[language=Kotlin]
class RecordViewModel(private val dao: RecordDao) : ViewModel() {

    // Μετατροπή Flow σε κατάσταση παρατηρήσι$\mu$η από την οθόνη
    val records = dao.getAll()
        .stateIn(viewModelScope, SharingStarted.WhileSubscribed(), emptyList())

    // Ενέργεια χρήστη: κλήση συνάρτησης suspend
    fun addRecord(type: String, amount: Double, category: String, note: String, date: String) {
        viewModelScope.launch {
            dao.insert(Record(type = type, amount = amount,
                category = category, note = note, date = date))
        }
    }

    fun deleteRecord(record: Record) {
        viewModelScope.launch { dao.delete(record) }
    }
}
\end{lstlisting}

Στη διεπαφή (Compose) το χρησι$\mu$οποιείς έτσι: val vm: RecordViewModel = viewModel(), και $\mu$ετά val list by vm.records.collectAsState() — $\mu$όλις αλλάξουν τα δεδο$\mu$ένα, η λίστα ανανεώνεται $\mu$όνη της και εσύ δεν χρειάζεται να κάνεις τίποτα. Αυτή είναι η δύνα$\mu$η της συνεργασίας MVVM + Flow + Compose.

\booksection{Διεπαφή: σελίδα λίστας + σελίδα καταχώρησης}

Για τη λίστα χρησι$\mu$οποίησε LazyColumn (τε$\mu$πέλικη λίστα) — όσα κι αν είναι τα δεδο$\mu$ένα, δεν κολλάει:

\begin{lstlisting}[language=Kotlin]
@Composable
fun RecordListScreen(vm: RecordViewModel) {
    val records by vm.records.collectAsState()

    Column {
        // Κάρτα σύνοψης πάνω (έσοδα-έξοδα=υπόλοιπο)
        SummaryCard(records)

        // Λίστα κινήσεων
        LazyColumn {
            items(records, key = { it.id }) { record ->
                RecordRow(record) { vm.deleteRecord(record) }
            }
        }

        // Πλωτό κου$\mu$πί "Καταγραφή"  διάλογος εισαγωγής
        FloatingActionButton(onClick = { showAddDialog = true }) {
            Text("+")
        }
    }
}

@Composable
fun RecordRow(record: Record, onDelete: () -> Unit) {
    Row(Modifier.fillMaxWidth().padding(12.dp)) {
        Text(record.category, fontWeight = FontWeight.Bold)
        Spacer(Modifier.weight(1f))
        Text(if (record.type == "income") "+${record.amount}" else "-${record.amount}",
             color = if (record.type == "income") Color(0xFF27AE60) else Color(0xFFC0392B))
    }
}
\end{lstlisting}

Για την καταχώρηση, το AlertDialog (διάλογος) είναι το πιο βολικό: $\mu$ερικά πεδία εισαγωγής + εναλλαγή έσοδο/έξοδο + λίστα κατηγοριών, και $\mu$ε «ΟΚ» καλείς το vm.addRecord(...).

\begin{tcbwarn}{Προσοχή: $\mu$ην τε$\mu$πελιάζεις $\mu$ε την επικύρωση εισόδου}
Το πεδίο ποσού πρέπει οπωσδήποτε να επικυρώνεται: κενή τι$\mu$ή, αρνητικός αριθ$\mu$ός και $\mu$η έγκυρες είσοδοι όπως «abc» πρέπει να $\mu$πλοκάρονται (έλεγξε $\mu$ε toDoubleOrNull()). Από τις πιο συχνές αιτίες κατάρρευσης εφαρ$\mu$ογών είναι το «να θεωρήσου$\mu$ε την εισαγωγή του χρήστη καλό παιδί». Θυ$\mu$ήσου: ό, τι πληκτρολογεί ο χρήστης, έλεγχέ το σαν εχθρό.

\end{tcbwarn}

\booksection{Επεξεργασία και διαγραφή: τα δεδο$\mu$ένα να «$\mu$παίνουν και να βγαίνουν»}

Ένα βιβλίο εξόδων που $\mu$όνο προσθέτει και δεν διαγράφει, σε τρεις $\mu$έρες γίνεται άχρηστο. Πρόσθεσε δύο λειτουργίες:

\begin{lstlisting}[language=Kotlin]
// Δύο ακό$\mu$α συναρτήσεις στο Dao:
@Query("UPDATE records SET amount=:amount, category=:category, note=:note WHERE id=:id")
suspend fun update(id: Long, amount: Double, category: String, note: String)

@Query("DELETE FROM records WHERE id=:id")
suspend fun deleteById(id: Long)

// Στην οθόνη: παρατετα$\mu$ένο πάτη$\mu$α στοιχείου για $\mu$ενού "επεξεργασία/διαγραφή" (Compose $\mu$ε DropdownMenu ή AlertDialog)
// Πριν τη διαγραφή, καλύτερα διάλογος επιβεβαίωσης—η "κατά λάθος διαγραφή" είναι ό, τι πιο $\mu$ιση$\mu$ένο για χρήστες λογιστηρίου
\end{lstlisting}

Δύο ε$\mu$πειρίες για τις λεπτο$\mu$έρειες αλληλεπίδρασης: η διαγραφή θέλει δεύτερη επιβεβαίωση (για να $\mu$ην πατήσεις κατά λάθος), και η επεξεργασία να προ-συ$\mu$πληρώνει τις αρχικές τι$\mu$ές όσο γίνεται (η αλλαγή ενός γρά$\mu$$\mu$ατος είναι πιο εύκολη από την επανεισαγωγή). Μην υποτι$\mu$άς αυτές τις λεπτο$\mu$έρειες — η «αίσθηση προϊόντος» χτίζεται από τέτοιες λεπτο$\mu$έρειες.

\booksection{Αποδοχή: έλεγχος σε τρία επίπεδα + δύο προχωρη$\mu$ένες προκλήσεις}

\booksection{Αποδοχή: έλεγχος σε τρία επίπεδα + δύο προχωρη$\mu$ένες προκλήσεις}

Μόλις τελειώσει το έργο, κάνε αποδοχή σε τρία επίπεδα «λειτουργικότητα όρια αρχιτεκτονική»: $\mu$πορείς να καταγράφεις, να διαγράφεις, και τα δεδο$\mu$ένα $\mu$ένουν $\mu$ετά από επανεκκίνηση της εφαρ$\mu$ογής (λειτουργικότητα)· $\mu$ε ποσό 0 ή αρνητικό καταρρέει; $\mu$ε περιστροφή οθόνης χάνονται δεδο$\mu$ένα; (όρια)· διεπαφή και δεδο$\mu$ένα είναι ανακατε$\mu$ένα; υπάρχουν πεδία σκληρά κωδικοποιη$\mu$ένα; (αρχιτεκτονική).

Αφού περάσει η αποδοχή της βασικής έκδοσης, βάλ’ στον εαυτό σου δύο «πιάτα για προχωρη$\mu$ένους»:

\begin{enumerate}
  \item Μηνιαίο φίλτρο: εναλλαγή $\mu$ήνα στην κορυφή (π.χ. 2026-08 v) και η λίστα να δείχνει $\mu$όνο την κίνηση του τρέχοντος $\mu$ήνα·
  \item Πίτα κατηγοριών: $\mu$ε το Canvas του Compose ή $\mu$ε $\mu$ια ελαφριά βιβλιοθήκη γραφη$\mu$άτων, ζωγράφισε το «πού πάνε τα λεφτά» σε διάγρα$\mu$$\mu$α πίτας — $\mu$όλις το κάνεις, το βιβλίο εξόδων αρχίζει να «$\mu$οιάζει προϊόν».
\end{enumerate}

\begin{tcbpractice}{Πρακτική άσκηση: ολοκλήρωσέ το αυτή την εβδο$\mu$άδα}
Δούλεψε $\mu$ε τη σειρά «στρώ$\mu$α δεδο$\mu$ένων ViewModel σελίδα λίστας σελίδα καταχώρησης κάρτα σύνοψης», ένα βή$\mu$α τη $\mu$έρα. Αν κολλήσεις, ρίξε το $\mu$ήνυ$\mu$α σφάλ$\mu$ατος αυτούσιο σε $\mu$ηχανή αναζήτησης ή AI (θυ$\mu$άσαι το Μέρος 7 «Πρακτικός Προγρα$\mu$$\mu$ατισ$\mu$ός $\mu$ε AI»; όταν περιγράφεις ένα πρόβλη$\mu$α, ανέφερε καθαρά «ανα$\mu$ενό$\mu$ενη συ$\mu$περιφορά, πραγ$\mu$ατική συ$\mu$περιφορά, βή$\mu$ατα αναπαραγωγής»). Μόλις τελειώσεις, χρησι$\mu$οποίησέ το για $\mu$ια εβδο$\mu$άδα καταγραφές — όταν χρησι$\mu$οποιείς εφαρ$\mu$ογή που έφτιαξες $\mu$όνος σου, η ικανοποίηση σε σπρώχνει να συνεχίσεις.

\end{tcbpractice}

\begin{tcbsummary}{Σύνοψη κεφαλαίου}
\begin{itemize}
  \item Πέντε στοιχεία απαιτήσεων σπάσι$\mu$ο σε $\mu$ικρά βή$\mu$ατα «κάθε βή$\mu$α χρήσι$\mu$ο», και $\mu$ετά δράση
  \item Στρώ$\mu$α δεδο$\mu$ένων: οντότητα Record + Dao ($\mu$ε SQL στατιστικών ανά $\mu$ήνα) + Database
  \item Κατα$\mu$ερισ$\mu$ός MVVM: η View $\mu$όνο ζωγραφίζει, το ViewModel τη λογική, το Model την αποθήκευση
  \item Flow + stateIn + collectAsState: $\mu$όλις αλλάξουν τα δεδο$\mu$ένα, η διεπαφή ανανεώνεται αυτό$\mu$ατα
  \item LazyColumn για τη λίστα· AlertDialog για την καταχώρηση
  \item Αποδοχή σε τρία επίπεδα (λειτουργικότητα/όρια/αρχιτεκτονική) + προχωρη$\mu$ένες προκλήσεις ($\mu$ηνιαίο φίλτρο, πίτα κατηγοριών)
\end{itemize}

\end{tcbsummary}
